\documentclass[11pt]{article}
\usepackage[T1]{fontenc}
\usepackage{lmodern}
\usepackage[margin=1in,top=.8in,bottom=.85in]{geometry}
\usepackage{amsmath,amssymb,amsthm,booktabs,microtype}
\usepackage[hidelinks]{hyperref}
\newtheorem{theorem}{Theorem}
\newtheorem{lemma}[theorem]{Lemma}
\newtheorem{corollary}[theorem]{Corollary}
\newcommand{\F}{\mathcal F}
\newcommand{\rev}{\mathrm{rev}}
\title{Factor frequencies and fifth powers in the base-$3/2$ Thue--Morse words}
\author{Blake Stone mathematics repository\\[-2pt]\small Proof developed with OpenAI Codex}
\date{November 26, 2026}
\begin{document}
\maketitle
\begin{abstract}
We prove Conjectures 7--10 in the journal version of Dekking's paper on the base-$3/2$
Thue--Morse sequence. For both the normalized digit-sum word and Dekking's variant, every
finite factor has a computable rational frequency, invariant under complementation and
reversal. The frequency is the same on every arithmetic progression of power-of-two step.
The analytic input is the dyadic single-letter theorem of Cassaigne, Espinoza, Rigo and
Stipulanti. Separately, an elementary descent proves that the critical exponent of each
word is exactly five. The only fifth powers in Dekking's word are $(010)^5$ and $(101)^5$,
both of frequency $1/180$.
\end{abstract}

\section{The two words and the questions}
Use zero-based positions. The normalized base-$3/2$ digit-sum parity word $t$ begins
$001110111110\cdots$ and satisfies
\begin{equation}
 t_0=0,\qquad t_{3n}=t_{3n+1}=t_{2n},\qquad t_{3n+2}=1-t_{2n+1}.\label{eq:rule}
\end{equation}
Dekking's word $x$, OEIS \href{https://oeis.org/A357448}{A357448}, begins
$010010101101\cdots$. Its rule sends a pair $ab$ to $a\bar a a$, where $\bar a=1-a$.
Write $\varphi(a)=a\bar a a$. Then $x=\varphi(t)$ \cite[Lemma 4]{cers}. For a direct check,
set $x_{3n}=x_{3n+2}=t_n$, $x_{3n+1}=\bar t_n$. Equation~\eqref{eq:rule} in the three
residue classes of $n$ gives $x_{2n}=t_n$, as required by the pair rule. With first bit
zero these equations determine the word.

Dekking's journal Conjectures 7, 8 and 9 concern existence of all factor frequencies,
invariance of those frequencies under bit-wise complement, and invariance under reversal,
respectively. Conjecture 7 predicts $\mu_x(00)=1/10$ and $\mu_x(01)=4/10$.
Conjecture 10, attributed to Jeffrey Shallit, asks whether the critical exponent is five
\cite[pp.~3--4]{dek}. The earlier arXiv version numbers these four statements 6--9.
Cassaigne, Espinoza, Rigo and Stipulanti proved uniform recurrence, symmetries of the
factor languages and single-letter frequencies. Their concluding remarks also predict
the two-letter frequencies for $t$ \cite[Sections 1.5 and 4]{cers}.

\begin{theorem}\label{main}
For $u=t$ or $x$, every binary word $w$ has a rational frequency $\mu_u(w)$. More precisely,
for every $h=2^a$ with $a\ge0$ and $b\ge0$,
\[
 \lim_{N\to\infty}\frac1N\#\{0\le j<N:u[hj+b,hj+b+|w|)=w\}=\mu_u(w).
\]
This value equals $\mu_u(\bar w)$ and $\mu_u(w^{\rev})$. Both words have critical exponent
five. Every period-$p$ factor of length at least $5p$ has exactly that length. For $t$ its
period is $p=1$. For $x$ it has $p=3$ and begins at a multiple of three.
\end{theorem}

\section{Decoding a finite factor}
Put $P(n)=\lfloor2n/3\rfloor$ and $e(n)={\bf1}_{n\equiv2\pmod3}$. With $\oplus$ denoting
addition modulo two, equation~\eqref{eq:rule} says $t_n=t_{P(n)}\oplus e(n)$. Fix a
nonempty word $w=w_0\cdots w_{m-1}$ and phase $r\in\{0,1,2\}$. Place the bit
$w_i\oplus e(r+i)$ in slot $P(r+i)-P(r)$ for every $i$. Two inconsistent demands give
$D_r(w)=\bot$. Otherwise the slots, read in order, form a word $D_r(w)$: $P$ advances by
zero or one, so there are no gaps. At a start $3k+r$ the factor is $w$ precisely when
$D_r(w)$ occurs at $2k+P(r)$. The symbol $\bot$ has zero indicator and zero frequency.

For $x$, use instead $P_x(n)=\lfloor n/3\rfloor$ and $e_x(n)={\bf1}_{n\equiv1\pmod3}$.
Call this decoder $E_r(w)$. Its parent word is read in $t$ at $k$.
For example $D_0(01)=\bot$, $D_1(01)=00$, and $D_2(01)=11$.

\begin{lemma}\label{length}
Every compatible $D_r(w)$ has length at most $m$. If $m\ge2$, phase zero has length
strictly less than $m$ whenever it is compatible. For $m\ge4$ all compatible phases
strictly shorten the word.
\end{lemma}
\begin{proof}
The length is $1+P(r+m-1)-P(r)$. Phase zero merges the first two positions.
Any interval of four consecutive positions includes a pair $3j,3j+1$, with the
same parent. This proves the assertions.
\end{proof}

We first define rational numbers $F(w)$ independently of limiting counts. Set $F(\bot)=0$,
$F(\varepsilon)=1$, and use the following base table for lengths one through three.
\begin{center}
\begin{tabular}{llc}
\toprule
Length & Words & $F(w)$\\
\midrule
1 & $0,1$ & $1/2$\\
2 & $00,11$ & $3/10$\\
  & $01,10$ & $1/5$\\
3 & $010,101$ & $1/20$\\
  & the other six words & $3/20$\\
\bottomrule
\end{tabular}
\end{center}
For $m\ge4$ the following identity defines $F$ by induction:
\begin{equation}
                3F(w)=F(D_0(w))+F(D_1(w))+F(D_2(w)).\label{eq:transfer}
\end{equation}
The same identity holds for $m=2,3$. For clarity the table below lists the decoding for
representatives under complement. Complement commutes with all three decoders.
\begin{center}
\begin{tabular}{rccc@{\qquad}rccc}
\toprule
$w$ & $D_0$ & $D_1$ & $D_2$ &$w$ & $D_0$ & $D_1$ & $D_2$\\
\midrule
$00$ & $0$ & $01$ & $10$ & $01$ & $\bot$ & $00$ & $11$\\
$000$ & $01$ & $010$ & $10$ & $001$ & $00$ & $011$ & $\bot$\\
$010$ & $\bot$ & $000$ & $\bot$ & $011$ & $\bot$ & $001$ & $11$\\
\bottomrule
\end{tabular}
\end{center}
Substitution of the base values proves the identity for these lengths.
Convexity gives $0\le F(w)\le1$ for all words.

\section{Why the limiting frequencies satisfy the calculation}
The only analytic input from \cite[Theorem 18]{cers} is that bits of $t$ on every progression
$2^a j+b$ have frequency $1/2$ each, with the denominator measured in progression terms.
The cited theorem uses the full prefix length as denominator and gives $2^{-a-1}$ for one
residue; dividing by the density $2^{-a}$ gives this formulation. An offset outside the
first residue period changes finitely many initial terms.

Let $A_{a,b}(w,N)$ be the relative count in Theorem~\ref{main} for $t$. Split the
progression index into $j=3k+d$, $0\le d<3$. Write $2^a d+b=3q_d+r_d$. The original
start is $3(2^ak+q_d)+r_d$. Its decoded factor is read at
$2^{a+1}k+2q_d+P(r_d)$. Since $2^a$ is coprime to three, $r_d$ runs once over $0,1,2$.
For $N=3K$ this gives an exact average of three parent counts at $K$.
For arbitrary $N$ the incomplete final triple changes normalized counts by $O(1/N)$.

We prove convergence to $F(w)$ by strong induction on $m$. The cases $m=0,1$ are already
known. At $m\ge2$, take
\[
 B_m=\sup_{a,b,\;|w|=m}\limsup_{N\to\infty}|A_{a,b}(w,N)-F(w)|.
\]
The supremum is taken \emph{after} the individual upper limits, and $0\le B_m\le1$.
Subtracting \eqref{eq:transfer} from the counting identity bounds the upper limit for
each given progression and word by the average of three upper limits. An inconsistent
branch has zero error. A shorter parent has zero error by induction. Phase zero is one
of these two cases. The other two errors are at most $B_m$, including when their dyadic
step doubles. Hence $B_m\le2B_m/3$, so $B_m=0$. It follows that $\mu_t(w)=F(w)$.
Applying the three phases of the $x$ decoder gives
\begin{equation}
                 \mu_x(w)=\tfrac13\sum_{r=0}^2 F(E_r(w)).\label{eq:sibling}
\end{equation}
For starts in a dyadic progression its parent step is $2^a$, so the same result holds
there. This proves existence and supplies an exact algorithm.

Complementing all demands preserves consistency and complements every resulting parent bit.
For reversal in the $t$ decoder, reflect positions by $n\mapsto1-n$. This exchanges
the two positions in each equal-parent pair and preserves the singleton positions. More
explicitly for $n\in\mathbb Z$,
$P(1-n)=-P(n)$ and $e(1-n)=e(n)$. Shifting a starting phase by $3j$ adds $2j$ to every
parent index and leaves slot differences unchanged. Thus reversing a length-$m$ word
permutes phases by $r\mapsto1-r-(m-1)\pmod3$ and reverses the decoded word, preserving
incompatibility. For $x$ use reflection $n\mapsto2-n$: its phase-1 sign is preserved and
$P_x(2-n)=-P_x(n)$. The base table is invariant under both operations. Induction in
\eqref{eq:transfer} for $m\ge4$, followed by \eqref{eq:sibling}, proves the frequency
symmetries. In particular,
\[
 \begin{aligned}\mu_t(00)=\mu_t(11)&=3/10, & \mu_t(01)=\mu_t(10)&=1/5,\\
 \mu_x(00)=\mu_x(11)&=1/10, & \mu_x(01)=\mu_x(10)&=2/5.\end{aligned}
\]

\section{How long can a repetition be?}
A factor $u[s,s+L)$ has period $p\ge1$ if its letters at distances $p$ agree wherever
both lie in the factor. Its exponent is its length divided by its least positive period.
The critical exponent is the supremum over all nonempty factors.

\begin{lemma}\label{edges}
In $t$ adjacent letters agree at every edge $3j$ and differ at every edge
$9j+1$ and $9j+5$. In $x$ they differ at every edge $3j$ and agree at every
edge $9j+2$. An edge $i$ compares positions $i,i+1$.
\end{lemma}
\begin{proof}
Equation~\eqref{eq:rule} gives the equalities at $3j$. At $9j+1,9j+2$ the two bits
are $t_{4j},\bar t_{4j}$. At $9j+5,9j+6$ they are $\bar t_{4j+2},t_{4j+2}$.
In $x$, each image $a\bar a a$ supplies a differing first edge. Its last bit at $9j+2$
and the first bit of the next image are $t_{3j},t_{3j+1}$, which agree.
\end{proof}

Suppose a period-$p$ factor of length $5p$ starts at $s$ and $3\nmid p$. Its edge
indicators are periodic with period $p$ within the available $5p-1$ edges. For every
offset $0\le i<p$, one of $s+i,s+i+p,s+i+2p$ is divisible by three, and all three
edges lie in the factor. Hence in $t$ every edge of the factor is equal, and in $x$
every edge is unequal. A constant run in $t$ has length at most five because every five
consecutive edge positions include $1$ or $5$ modulo nine. An alternating run in $x$
has length at most nine because every nine consecutive edges include $2$ modulo nine.
It follows that in $t$ only $p=1$ is possible. In $x$ every $p\ge2$ is excluded, while
$p=1$ contradicts the required alternation.

\begin{lemma}\label{pull}
If $t[s,s+L)$ has period $3q<L$, the contiguous parent interval
$[P(s),P(s+L-1)]$ has period $2q$. If $x[s,s+L)$ has period $3q<L$, its contiguous
parent interval $[P_x(s),P_x(s+L-1)]$ in $t$ has period $q$.
\end{lemma}
\begin{proof}
For $t$, the signs agree at a separation of $3q$, and $P(n+3q)=P(n)+2q$.
The indices $s,\ldots,s+L-1-3q$ visit every parent from $P(s)$ through
$P(s+L-1)-2q$ because $P$ takes steps of size zero or one. Each required comparison
is therefore inherited. Use $P_x(n+3q)=P_x(n)+q$ for the second assertion.
\end{proof}

A length-$15q$ factor in $t$ covers at least $10q$ parent positions. If it had period
$3q$, Lemma~\ref{pull} would produce a fifth power of period $2q$. Strong induction
on the period, starting with the $3\nmid p$ case, now shows that any length-$5p$
periodic factor in $t$ must have $p=1$: in the divisible case one would need $2q=1$.
No such length can be exceeded, since a longer factor would contain its first $5p$
letters and a constant run cannot exceed five.

A length-$15q$ factor of period $3q$ in $x$ covers at least $5q$ positions in $t$
with period $q$. Consequently $q=1$. A length-$15$ window starting inside a triple
covers six parent positions and is impossible, and a length-$16$ window from any
alignment also covers at least six. The fifth powers in $x$ can therefore only be
aligned images of five identical bits. Both occur: $t[6,11)=11111$ and
$t[87,92)=00000$, by \eqref{eq:rule}. Their images have least period three.
This proves the repetition assertions of Theorem~\ref{main}, including attainment of
the critical exponent. Finally \eqref{eq:transfer} gives
$F(00000)=F(11111)=1/60$. Each aligned image has one possible phase out of three,
so each fifth power has frequency $1/180$ in $x$.

\paragraph{Formal verification.}
The accompanying Lean development proves the period bounds directly from the two word
definitions. For frequencies it takes the full-prefix statement of Theorem~18 of
\cite{cers} as an explicit hypothesis and derives the limits, symmetries and displayed
values. The repository also includes an exact rational calculator for any finite word.

\begin{thebibliography}{9}
\bibitem{dek} F. M. Dekking, \emph{The Thue--Morse sequence in base $3/2$},
Journal of Integer Sequences \textbf{26} (2023), Article 23.2.3.
\url{https://cs.uwaterloo.ca/journals/JIS/VOL26/Dekking/dek25.pdf}.
\bibitem{cers} J. Cassaigne, B. Espinoza, M. Rigo, and M. Stipulanti,
\emph{Symbols frequencies in the Thue--Morse word in base $3/2$ and related conjectures},
arXiv:2602.21895 (2026). \url{https://arxiv.org/abs/2602.21895}.
\end{thebibliography}
\end{document}
